%In the current paper we introduce \emph{timed scenario expressions}
%and solve the problem of realisability of sets of
%scenarios as scenario expressions.
%That is, given a global model of the behaviours of a system, captured
%by a set of scenarios, $\mathcal{S}$, we produce a scenario expression that expresses
%the same set of behaviours that are allowed by the members of $\mathcal{S}$.
%So our contribution is two fold:
We introduced the notion of timed scenario expressions (expressions for
short) and their semantics.
A distributed timed scenario (DTS) is an expression.
Expressions are also built by applying three operators, $|$, $\circ$ and $||$, to
expressions. The operator
$|$ specifies a choice between two alternatives, while $\circ$ and
$||$ indicate the concatenation and the parallel composition of two
expressions.
%, respectively.
We defined $\Sem{E}$, the semantics of expression $E$,
in terms of the set of behaviours that are
specified by $E$, and defined the consistency of $E$ in
terms of its semantics.

%A behaviour \Behaviour{B} of $\Xi$ is obtained by merging
%the behaviours of the individual components in $\Xi$ in such a way that
%the order of events does not violate $\ParSE{\DScenario}$, the
%underlying (extended) partial order on events of \DScenario.
%Moreover, each event that occurs in more than one component (called a
%synchronizing event) occurs only once.

%We formally define the consistency of distributed scenarios, in terms
%of their semantics, and formulate a criterion
%(Theorem \ref{thm-consistency}), as well as a practical method
%for determining the consistency of a DTS.

We tackled the problem of realisability of sets of (sequential)
scenarios as scenario expressions.
We developed an algorithm that, given a set of scenarios,
$\mathcal{S}$, produces an expression $E$ that realises
$\mathcal{S}$, i.e, $\Sem{\mathcal{S}}=\Sem{E}$.

The set $\mathcal{S}$ can sometimes be realised by a single
DTS. But, in general, realising $\mathcal{S}$
involves dividing it into partitions, in such a way that each
partition can be realised by a single DTS. The
choice expression whose alternatives are these distributed
scenarios would realise $\mathcal{S}$.

%Our goal is to produce an expression that includes a reasonably few
%number of alternatives, in particular, we avoid building a
%trivial solution in which every member of $\mathcal{S}$ is a separate
%alternative, when other solutions exist.

We investigated the fundamental reasons when $\mathcal{S}$ is not
realisble by a DTS.
%, based on which we
This investigation resulted in developing
two methods for effectively dividing $\mathcal{S}$.
One is based on reasoning about the underlying partial order in
$\mathcal{S}$, and the other is based on the time constraints of
$\mathcal{S}$.
%We introduce two different measures, K-ranking and c-ranking, for
%evaluating various divisions of $\mathcal{S}$.
%Intuitively, the K-ranking and c-ranking of a division, $D$,
%are measures of similarity of the
%sequences of events, and of the time constraints
%between the scenarios in each partition of $D$.
%We call the divisions that are obtained based on this analysis
%order-based divisions (OBDs) and constraint-based divisions (CBDs).
%A division is of the form $D=(\mathcal{S}_1,\mathcal{S}_2)$, where
%$\mathcal{S}_1\cup\mathcal{S}_2=\mathcal{S}$.
We identified the necessary and sufficient conditions for realisability
of a set of scenarios as a single DTS (Theorem \ref{thm-realisable}),
and thereby solved the problem that was left open in our earlier
work \cite{DTS-2025}.

To illustrate the effectiveness of our approach
we compared our algorithm with one that performs a complete
breadth-first search to find
%a division that results in
an optimal
solution, i.e., one with the minimum number of alternatives.
Our experimental results confirm the effectiveness and efficiency of
our algorithm.
%While we do not claim optimality, in all our experiments our algorithm
%always returned an optimal solution.

%The algorithms for determining consistency of expressions and
%realisability of sets of scenarios as expressions are
%implemented and available on request.

As our future work we plan to incorporate the Kleene star
operator in scenario expressions, and to explore the usefulness of machine learning for
solving the problem of realisability.
